\section*{Sources and editorial status}
The \href{https://protobuf.dev/programming-guides/encoding/}{Protocol Buffers encoding
guide} describes seven-bit unsigned varints and the ZigZag signed mapping. This exercise
uses those arithmetic ideas without field tags or Protobuf message framing. Its exact
delta width, capacities, status values and canonicality rules are the local contract.
The \href{https://webassembly.github.io/spec/core/binary/values.html\#integers}{WebAssembly
integer encoding specification} permits some longer equivalent LEB128 encodings within
its width limits. This project's rejection of nonminimal terminal zero payloads must
therefore be read as an explicit project policy.

Instruction behavior is checked against MARS revision \texttt{f1ddaef8}, especially
\texttt{WideInstructionCatalog.kt} for doubleword arithmetic and shifts, and the native
signed comparisons. The executable prototypes are checked against an independent
arithmetic model, including malformed input and capacity priorities. Testing a paired
encoder and decoder alone would not prove agreement with the format.

This directly edited native chapter now contains sixty steps, fifty cumulative edits,
five supplied frame/control excerpts, sixty quizzes, one hundred and twenty progressive
hints and eighty-four exact observation recipes. The construction checks cover both an
isolated function at offset addresses and the actual program caller. They distinguish
provisional construction states from the completed item-atomic interface.

Student and instructor archives, public per-step check packaging, enforced progression
and actual light-theme IDE captures remain pending. The four capture slots describe
planned observations; they are not screenshots. This chapter is not yet included in
the main manuscript's twenty-two text packages. Human review and independent
accessibility validation remain pending; the full thirty-five-project objective is active.
